% Generated draft: requires mathematical and source review.
\documentclass[11pt]{article}
\usepackage[margin=25mm]{geometry}
\usepackage{fontspec}
\setmainfont{Latin Modern Roman}
\usepackage{amsmath,amssymb,mathtools,unicode-math}
\setmathfont{Latin Modern Math}
\usepackage{enumitem,needspace,xurl,hyperref}
\setlength{\parindent}{0pt}
\setlength{\parskip}{5pt}
\setlength{\emergencystretch}{3em}
\newcommand{\N}{\mathbb{N}}
\newcommand{\Z}{\mathbb{Z}}
\newcommand{\Q}{\mathbb{Q}}
\newcommand{\R}{\mathbb{R}}
\newcommand{\C}{\mathbb{C}}
\newcommand{\F}{\mathbb{F}}
\newcommand{\A}{\mathbb{A}}
\newcommand{\PP}{\mathbb{P}}
\newcommand{\E}{\mathbb{E}}
\newcommand{\eps}{\varepsilon}
\newcommand{\dd}{\,\mathrm{d}}
\begin{document}
\section*{UnsolvedMath ID: 30004685}
\textbf{Source problem code:} catalogue entry
\subsection{Definitions and mathematical statement}
The supplied record poses a research programme rather than a single quantified conjecture. The following notation explains its terms without selecting hypotheses absent from the source.

Let \(X\) denote an unspecified state space, \(T\) an unspecified time set, and \(\mathcal{A}\subseteq X^T\) a class of admissible trajectories of a dynamical system. A sensor or measurement map is a map \(h_i:X\to Y_i\), where \(i\) indexes sensors or measurement modalities and \(Y_i\) is its output space. For a collection \(I\) of sensor--time pairs, define the ideal observation map by
\[
\mathcal{O}_I:\mathcal{A}\longrightarrow\prod_{(i,t)\in I}Y_i,
\qquad
\mathcal{O}_I(x)=\bigl(h_i(x(t))\bigr)_{(i,t)\in I}.
\]
Here sensor placement concerns the choice of available sensors or their locations; measurement-map discovery concerns identifying or choosing the maps \(h_i\); time-delay methods use observations at multiple appropriately spaced times; and multimodal fusion combines outputs from different measurement modalities. The source specifies neither the admissible choices nor their constraints.

For a specified target time \(t_*\in T\), exact noiseless full-state reconstruction would mean finding a map
\[
R:\mathcal{O}_I(\mathcal{A})\to X
\quad\text{such that}\quad
R(\mathcal{O}_I(x))=x(t_*)
\quad\text{for every }x\in\mathcal{A}.
\]
This is a definition of one possible reconstruction target, not an assertion that such a map exists or that the source intends this particular target. Reduced coordinates may be represented abstractly by a map \(r:X\to Z\), but the source supplies no dimension, regularity, dynamical fidelity, or other criterion making \(Z\) reduced or the coordinates useful. The background also includes noisy measurements; it gives no noise model, error metric, stability requirement, or probabilistic quantifier, so the displayed noiseless definition does not formalize that additional requirement.

\textbf{Source question.} What are the mathematical limits and possibilities of sensor placement, measurement-map discovery, time-delay methods, and multimodal data fusion for reconstructing the full state of a complex dynamical system and producing useful reduced coordinates?

A precise theorem or conjecture cannot be extracted uniquely without specifying the dynamical-system class, available information, observation constraints, reconstruction target, noise assumptions, and meaning of useful reduced coordinates. These qualifications remain unresolved.
\subsection{Short English statement}
Determine when sensor placement, learned measurement maps, delayed observations, and multimodal measurements permit full-state reconstruction and useful reduced coordinates. The supplied question leaves the system class, noise assumptions, admissible methods, and success criteria unspecified, so it remains a research programme requiring further formulation.
\subsection{Sources}
\begin{itemize}
\item Local UnsolvedMath registry: supplied complete record for permanent ID \(30004685\), titled \emph{State Reconstruction and Reduced Coordinates from Sparse Sensors}, in the Oberwolfach Reports Open Problems collection; exact source code \texttt{catalogue entry}. No individual proposer is supplied.
\item Supplied report citation: \emph{Mini-Workshop: Analysis of Data-driven Optimal Control (2022)}, Oberwolfach Reports, pp.\ 1209--1257, DOI \texttt{10.4171/owr/2021/23}. \url{https://doi.org/10.4171/owr/2021/23}. The registry records proposed year 2021; its relationship to the citation's 2022 date is unresolved. No specific problem page or workshop contributor is identified.
\item Casdagli et al., state reconstruction with noise: \url{https://doi.org/10.1016/0167-2789(91)90222-U}.
\item Serpas et al., sensor location for nonlinear systems: \url{https://doi.org/10.1016/j.compchemeng.2012.07.014}.
\item Sendiña-Nadal and Letellier, nonlinear network reconstruction: \url{https://doi.org/10.1063/5.0090239}.
\item Supplied literature assessment, checked \texttt{2026-08-22}: the record labels the problem partially solved and reports restricted rigorous or constructive results involving Takens-type embeddings, observability-Gramian placement, and nonlinear network reconstruction, while describing general limits of full-state recovery from multimodal noisy measurements as open. This is supplied dated curation derived from the OWR record, not an independently verified assessment; its status and cited claims require independent verification.
\end{itemize}
\end{document}
